\chapter{Mapa de Dependencias entre Capítulos}
\label{chap:deps}

Este capítulo describe el grafo de prerrequisitos entre capítulos,
permitiendo al lector planificar una lectura no lineal cuando sea necesario.

\section{Base Universal}

La Parte I (Fundamentos) es el único punto de entrada verdadero. No tiene
dependencias previas y es requisito para todas las demás partes.

\section{Dependencias por Parte}

\begin{itemize}
\item \textbf{Parte II} (El Núcleo Lisp) depende de la Parte I: sintaxis,
  tipos, control de flujo.
\item \textbf{Parte III} (El Puente Python) depende de las Partes I y II:
  funciones, símbolos, macros.
\item \textbf{Parte IV} (Programación Híbrida) depende de las Partes I, II y
  III: puente Python, CLOS, macros.
\item \textbf{Parte V} (Meta-Programación) depende de las Partes I--IV.
\item \textbf{Parte VI} (Matemáticas Aplicadas) depende de las Partes I, III
  y IV: \texttt{numpy-fast}, \texttt{sympy-bridge}, \texttt{defmath}.
\item \textbf{Parte VII} (Ingeniería de Sistemas) depende de las Partes I y
  II: compilación, CLI, sistema.
\item \textbf{Parte VIII} (Estudios de Caso) depende de las Partes I--VI
  según el caso de estudio específico.
\end{itemize}

\section{Matriz de Prerrequisitos}

\begin{center}
\small
\begin{tabular}{lccccccccc}
\toprule
\textbf{Capítulo} & \textbf{I} & \textbf{II} & \textbf{III} & \textbf{IV} & \textbf{V} & \textbf{VI} & \textbf{VII} \\
\midrule
I (1.1--1.8) & -- & -- & -- & -- & -- & -- & -- \\
II (2.1--2.9) & $\bullet$ & -- & -- & -- & -- & -- & -- \\
III (3.1--3.9) & $\bullet$ & $\bullet$ & -- & -- & -- & -- & -- \\
IV (4.1--4.11) & $\bullet$ & $\bullet$ & $\bullet$ & -- & -- & -- & -- \\
V (5.1--5.8) & $\bullet$ & $\bullet$ & $\bullet$ & $\bullet$ & -- & -- & -- \\
VI (6.1--6.8) & $\bullet$ & -- & $\bullet$ & $\bullet$ & -- & -- & -- \\
VII (7.1--7.7) & $\bullet$ & $\bullet$ & -- & -- & -- & -- & -- \\
VIII (8.1--8.7) & $\bullet$ & $\bullet$ & $\bullet$ & $\bullet$ & $\circ$ & $\circ$ & $\circ$ \\
\bottomrule
\end{tabular}
\end{center}

Leyenda: $\bullet$ = requisito directo, $\circ$ = requisito parcial.

\section{Capítulos Autocontenidos}

Los siguientes capítulos pueden leerse de forma independiente, marcados con
$(*)$ para lectura no lineal:

\begin{itemize}
\item Capítulo 2.7 -- LOOP avanzado $(*)$
\item Capítulo 3.9 -- FFI directa $(*)$
\item Capítulo 4.8 -- Persistencia $(*)$
\item Capítulo 7.3 -- Gestión de paquetes pip $(*)$
\item Capítulo 7.5 -- Pruebas $(*)$
\end{itemize}

\section{Recomendación de Orden}

Para cada ruta de aprendizaje (ver Capítulo~\ref{chap:roadmap}, se recomienda
seguir el orden numérico dentro de las partes seleccionadas. Los capítulos
autocontenidos pueden intercalarse según el interés del lector.
